\zlabel{triptych:concise:themes:start}
\section{The Propers: Themes and Movement}

The Last Sunday after Pentecost gives thanks for deliverance already
received and asks for strength to endure until Christ gathers his people.
Those who trust God to hear them also ask him to change what they desire.

\subsection{Promised peace and fruitful wills}

The Introit begins with the Lord's promise to the exiles in Babylon.
Jeremiah's letter tells them to build, plant and pray for the city where they
live. Against the hope of a quick return encouraged by false prophets,
Jerome explains that the exiles will be gathered at the time God promised.
Jerome says that their prayer prepares them to receive the promised
return.\footnote{Jer 29:4--14;
Jerome, \work{In Ieremiam} V, \textit{Vers. 10 seqq.} on Jer 29:10--14a
and opening of \textit{Vers. 14 seqq.} on Jer 29:14b, PL 24,
cols. 893--894.}

The Introit's psalm verse speaks of captivity reversed. Yet the complete Psalm
84(85) still asks for conversion and renewed life. Augustine hears the
prophet announcing a future restoration as accomplished. He finds in
Jacob's captivity the bondage to sin that Paul describes in Romans.
The Collect calls its petitioners faithful and asks God to stir their
wills for fruitful work. Fromage explains that good works done by grace
prepare the faithful to receive further grace.\footnote{Augustine,
\work{Enarrationes in Psalmos}, Ps 84, Latin §§4, 6; Fromage,
continuation of \work{The Liturgical Year}, vol. XI, p. 503.}

\subsection{Thanksgiving and the cry from the depths}

The Epistle prays for people whose faith and love Paul already knows.
Knowledge of God's will is to produce good work, growth, endurance and
thanksgiving. Chrysostom explains that their existing grace is a reason for
further encouragement. He likens them to competitors cheered on when
they are near victory. The Father has rescued them from darkness and
brought them into the beloved Son's kingdom. They still need strength to
live there faithfully.\footnote{Col 1:9--14; Chrysostom,
\work{Homilies on Colossians} 2.}

The Gradual praises deliverance, although the complete Psalm 43(44)
describes faithful people suffering defeat. Augustine hears the future
rescue spoken of with prophetic certainty. Bellarmine hears thanksgiving
for rescues already received.

The Alleluia then asks to be heard from the depths. In the historical
Lasance/Walsh wording it sings ``Lord, hear my prayer.'' Augustine
distinguishes the person who cries out from the person who has become
content to remain in sin. Worshippers who thank God for rescue in the
Gradual can still call on him in their need.\footnote{Augustine,
\work{Enarrationes in Psalmos}, Ps 43, Latin §§8--9, NPNF Ps XLIV
§§6--7, and Ps 129 §§1--2; Bellarmine, \work{Commentary on the Book
of Psalms}, Psalm 43, native vv. 7--8 (Vulgate vv. 8--9);
Lasance/Walsh, Last Sunday after Pentecost, Alleluia, p. 738.}

\clearpage
\subsection{Flight, false assurances and the public gathering}

In the Gospel, Jesus warns people in Judea to flee without retrieving
possessions. The difficulties of pregnancy, nursing, winter and Sabbath
concern actual escape. Chrysostom explains them before distinguishing the
Judean catastrophe from Christ's final coming. The warning requires prudent
action, and Jesus tells the disciples to pray about the time of the flight.
He promises to shorten the days for the elect.\footnote{Mt 24:15--35; Chrysostom,
\work{Homilies on Matthew} 76.}

The later warnings concern false claimants who offer a Christ hidden in a
desert or an inner room. Jesus compares his own coming to lightning visible
across the sky and promises the angels' gathering of the elect.
Chrysostom interprets ``this generation'' as the believing people who
will survive the trials as a people. The promise concerns the Church's
endurance; it does not say that each disciple avoids suffering or death.

In the interpretation of deliverance begun, those who come through one danger
still await the gathering Christ promises. In the interpretation of converted
desire, the false claimants expose how the desire for reassurance can make
an impressive sign seem more trustworthy than Christ's words.\footnote{Chrysostom,
\work{Homilies on Matthew} 77, on Mt 24:33--35; compare Aquinas,
\work{Super Matthaeum} 24, lectio 3.}

\subsection{Offering, confident prayer and medicinal reception}

The Offertory antiphon repeats the opening cry from Psalm 129. In the
Secret the Church asks God to receive the people's offerings and prayers
and turn their hearts toward him. Fromage explains that they ask to pass
from earthly attachment to heavenly desire.\footnote{1962 Missal, p. 424,
nos. 1775--1776; Fromage, continuation of \work{The Liturgical Year},
vol. XI, p. 513.}

The Trinity Preface gives thanks directly to the Father. It confesses the
Father, Son and Holy Spirit as distinct persons of one divine essence and
equal majesty. The worshippers' praise with the angels can express the
heavenly desire for which they prayed in the Secret, the desire for communion
with God.\footnote{1962 Missal, p. 293, no. 1082; Lasance/Walsh,
Most Holy Trinity Preface and Sanctus, pp. 773--774.}

In the Communion Christ tells those who ask in prayer to believe they will
receive. Theophylact understands this faith as turning one's whole heart
toward God. In Mark's next verse Christ tells those standing in prayer to
forgive so that the Father will forgive them. Bede sees in those words an
assurance to those seeking pardon who first forgive the people who have
wronged them.\footnote{Mk 11:22--25;
Theophylact of Ohrid, \work{In Marcum}, Basel 1525, fols. 69r--v, on
11:24--25; Bede, \work{In Marci evangelium expositio} III, on 11:25,
registered Latin transcription, line 339.}

The Collect asks for further remedies. After Communion, the worshippers ask
the mysteries they have received to cure the faults that remain.\footnote{1962
Missal, p. 423, no. 1770, and p. 424, no. 1778; Fromage, continuation of
\work{The Liturgical Year}, vol. XI, pp. 503, 513.}
\zlabel{triptych:concise:themes:end}
\clearpage
